% problem_id: S046-thm:w-hinit
% source_id: S046 (HoTT Book)
% source_file: induction.tex
% statement_locator: induction.tex lines 524--526
% proof_locator: induction.tex lines 528--562
% env: thm
% pairing: tight (verified)
% status: complete (statement + author proof verbatim + reason + locators)
%
\begin{thm}\label{thm:w-hinit}
For any type $A : \type$ and type family $B : A \to \type$, the $\w$-algebra $(\wtype{x:A}B(x), \supp)$ is h-initial.
\end{thm}

\begin{proof}[Sketch of proof]
Suppose we have $A : \type$ and $B : A \to \type$,
and consider the associated polynomial functor $P(X)\defeq\sm{x:A}(B(x)\to X)$.
Let $W \defeq \wtype{x:A}B(x)$.  Then using
the $\w$-introduction rule from \cref{sec:w-types}, we have a structure map $s_W\defeq\supp: PW \rightarrow W$.
We want to show that the algebra $(W, s_W)$ is
h-initial. So, let us consider another algebra $(C,s_C)$ and show that the type $T\defeq \whom_{A,B}((W, s_W),(C,s_C)) $
of  $\w$-homomorphisms from $(W, s_W)$ to $(C, s_C)$ is contractible. To do
so, observe that the $\w$-elimination rule and the $\w$-computation
rule allow us to define a $\w$-homomorphism $(f, s_f) : (W, s_W) \rightarrow (C, s_C)$,
thus showing that $T$ is inhabited. It is furthermore necessary to show that for every $\w$-ho\-mo\-mor\-phism $(g, s_g) : (W, s_W) \rightarrow (C, s_C)$, there is an identity proof
\begin{equation}
\label{equ:prequired}
p :  (f,s_f) = (g,s_g).
\end{equation}
This uses the fact that, in general, a type of the form $(f,s_f) = (g,s_g) $
is  equivalent to the type of what we call \define{algebra $2$-cells}
\indexdef{algebra!2-cell}%
from $f$ to $g$, whose canonical
elements are pairs of the form $(e, s_e)$, where $e : f=g$ and $s_e$ is a higher identity proof between the identity proofs represented by the following pasting diagrams:
\[
\xymatrix{
PW \ar@/^1pc/[r]^{Pg}   \ar[d]_{s_W}   \ar@/_1pc/[r]_{Pf} \ar@{}[r]|{Pe}
& PC \ar[d]^{s_C}  \\
W  \ar@/_1pc/[r]_f  \ar@{}[r]^{s_f} & C } \qquad
\xymatrix{
PW \ar@/^1pc/[r]^{Pg}   \ar[d]_{s_W} \ar@{}[r]_(.52){s_g}  & PC \ar[d]^{s_C}  \\
W \ar@/^1pc/[r]^g  \ar@/_1pc/[r]_f  \ar@{}[r]|{e} & C }
\]
In light of this fact, to prove that there exists an element as in~\eqref{equ:prequired}, it is
sufficient to show that there is an algebra 2-cell $(e,s_e)$ from $f$ to $g$.
The identity proof $e : f=g$ is now constructed by function extensionality and
$\w$-elimination so as to guarantee the existence of the required identity
proof $s_e$.
\end{proof}
